\chapter{Relativistic Quantum Mechanics and Local Operator Structure}
\label{ch:relativistic-qm-local-operator-structure}
\ConventionsLorentzian


\section{Hilbert-Space Quantum Data}

\begin{definition}[Hilbert-space quantum datum]
\label{def:hilbert-space-quantum-datum}
Let \(\Hilb\) be a complex Hilbert space. A pure state is a ray in \(\Hilb\).
A mixed state is a positive trace-class operator \(\rho\) on \(\Hilb\)
satisfying
\[
  \rho\ge0,
  \qquad
  \operatorname{Tr}_{\Hilb}\rho=1.
\]
For a bounded operator \(A\in\mathcal B(\Hilb)\), the expectation value in
the state \(\rho\) is
\[
  \langle A\rangle_\rho=\operatorname{Tr}_{\Hilb}(\rho A).
\]
For an unbounded operator, the operator domain and the domain of every product
appearing in a matrix element are part of the data.
\end{definition}

\begin{definition}[Hamiltonian time evolution]
\label{def:hamiltonian-time-evolution}
A Hamiltonian time evolution is specified by a self-adjoint operator \(H\) and
the unitary one-parameter group
\[
  U(t)=\ee^{-\ii tH},
  \qquad
  A(t)=U(t)^{-1}AU(t).
\]
Stone's theorem gives the converse: every strongly continuous unitary
one-parameter group has a unique self-adjoint generator.  When \(H\) is first
given as a formal differential expression, choosing its self-adjoint
realization is part of the definition of the quantum system.
In a stable sector \(H\) is bounded below. Spacetime localization is additional
structure: a local relativistic theory assigns operators or algebras to
spacetime regions.
\end{definition}

\section{Rigged Hilbert Spaces and Continuous Spectra}
\label{sec:rigged-hilbert-continuous-spectrum}

The Hilbert space contains the normalizable states. The momentum kets, energy
kets, position kets, and plane-wave scattering states used throughout the
monograph are generalized vectors. Their precise home is a rigged Hilbert
space, also called a Gelfand triple.

\begin{definition}[Rigged Hilbert space]
\label{def:rigged-hilbert-space}
A rigged Hilbert space consists of a dense subspace \(\Phi\subset\Hilb\),
equipped with a locally convex topology finer than the Hilbert norm topology,
together with the continuous inclusions
\[
  \Phi\subset\Hilb\subset\Phi^\times .
\]
Here \(\Phi^\times\) is the continuous antilinear dual of \(\Phi\). The second
inclusion sends \(h\in\Hilb\) to the functional
\[
  \phi\longmapsto \langle \phi,h\rangle_{\Hilb},
  \qquad \phi\in\Phi .
\]
Thus every Hilbert vector defines a generalized vector, but
\(\Phi^\times\) is larger than \(\Hilb\). Distributional eigenvectors live in
\(\Phi^\times\).
\end{definition}

The basic model is
\[
  \mathcal S(\mathbb R^d)
  \subset
  L^2(\mathbb R^d,\dd^d x)
  \subset
  \mathcal S'(\mathbb R^d),
\]
where \(\mathcal S\) is the Schwartz space and \(\mathcal S'\) is the space
of tempered distributions. The position and momentum kets are the tempered
distributions defined by
\[
  \langle x|\psi\rangle=\psi(x),
  \qquad
  \langle p|\psi\rangle=\widetilde\psi(p),
  \qquad
  \psi\in\mathcal S(\mathbb R^d),
\]
with the Fourier convention chosen in the relevant chapter. The identities
\[
  \langle x|x'\rangle=\delta^{(d)}(x-x'),
  \qquad
  \langle p|p'\rangle=(2\pi)^d\delta^{(d)}(p-p'),
\]
and
\[
  \int\dd^d x\,|x\rangle\langle x|=\mathbf 1,
  \qquad
  \int\frac{\dd^d p}{(2\pi)^d}\,|p\rangle\langle p|=\mathbf 1
\]
are weak identities: after pairing with test vectors in \(\Phi\) on the left
and right, they become ordinary identities for functions or distributions.
Point evaluation and Fourier evaluation are continuous linear functionals on
\(\mathcal S(\mathbb R^d)\), hence define elements of
\(\mathcal S'(\mathbb R^d)\).  Pairing the displayed resolution of the
identity with \(\varphi,\psi\in\mathcal S(\mathbb R^d)\) gives respectively
\(\int \overline{\varphi(x)}\psi(x)\dd^dx\) and the same inner product written
after Fourier inversion.  The delta distributions are precisely the kernels
which reproduce these pairings, so the displayed bra-ket identities are kernel
statements, not equalities between Hilbert vectors.

\begin{definition}[Generalized eigenvector in a rigged Hilbert space]
\label{def:generalized-eigenvector-rigged-hilbert}
For a self-adjoint operator \(A\) with spectral measure \(E_A\), a number
\(\lambda\) is an eigenvalue in the Hilbert-space sense precisely when
\[
  E_A(\{\lambda\})\Hilb\ne0 .
\]
If \(E_A(\{\lambda\})=0\) but \(\lambda\) lies in the continuous spectrum,
there is no normalizable eigenvector with that eigenvalue. A generalized
eigenvector at \(\lambda\) is instead a functional
\[
  F_\lambda\in\Phi^\times
\]
such that, on a common invariant test domain,
\[
  F_\lambda(A\phi)=\lambda F_\lambda(\phi),
  \qquad
  \phi\in\Phi .
\]
This is the mathematical meaning of a plane wave of sharp momentum or an
energy eigenstate in a continuous spectrum. It is a distributional object
used to write kernels. Physical Hilbert vectors are obtained by smearing such
kernels with square-integrable wave packets.
\end{definition}

\begin{definition}[Direct-integral coordinate ket]
\label{def:direct-integral-coordinate-ket}
Equivalently, the spectral theorem may be written as a direct integral
\[
  \Hilb
  \simeq
  \int_X^\oplus \Hilb_\lambda\,\dd\nu(\lambda),
  \qquad
  A
  =
  \int_X^\oplus \lambda\,\mathbf 1_{\Hilb_\lambda}\,\dd\nu(\lambda).
\]
The label \(\lambda\) indexes fibers over the spectrum. A generalized ket
\(|\lambda,a\rangle\) is not an element of \(\Hilb\).  It denotes the
distributional coordinate functional that extracts the \(a\)-component in the
fiber \(\Hilb_\lambda\) of the direct integral.  Delta-normalization records
the kernel of the identity relative to the measure \(\dd\nu\). Changing
\(\dd\nu\) rescales these generalized kets and therefore changes the displayed
delta factors; the Hilbert-space vector obtained after smearing is unchanged.
\end{definition}

\begin{definition}[Relativistic one-particle rigging]
\label{def:relativistic-one-particle-rigging}
For a relativistic one-particle mass shell this gives, for example,
\[
  \Phi_1
  =
  C_c^\infty(\Sigma_m^+;V)
  \subset
  \Hilb_1
  =
  L^2(\Sigma_m^+,\dd\mu_m;V)
  \subset
  \Phi_1^\times ,
\]
where \(V\) is the finite-dimensional little-group representation space. The
little-group label convention is fixed by choosing an orthonormal basis
\(\{e_\sigma\}\) of \(V\) and writing
\[
  \langle e_\tau,e_\sigma\rangle_V=\delta_{\tau\sigma}.
\]
Thus, in a bra-ket expression, the first Kronecker index is the label carried
by the bra and the second is the label carried by the ket; since the Kronecker
delta is symmetric, \(\delta_{\tau\sigma}=\delta_{\sigma\tau}\).  The symbol
\(|p,\sigma\rangle\) is an element of \(\Phi_1^\times\), and
\[
  \langle p,\tau|p',\sigma\rangle
  =
  (2\pi)^{D-1}2p^0\,
  \delta^{(D-1)}(\vec p-\vec p')
  \delta_{\tau\sigma}
\]
is the distribution kernel for the identity in the covariant mass-shell
measure. All scattering kernels written later with sharp external momenta
are to be read in this sense. The wave-packet formulation is the primary
Hilbert-space statement; the plane-wave formula is the corresponding
distributional kernel.
\end{definition}

\section{Poincare Symmetry}

\begin{definition}[Poincare-covariant vacuum Hilbert datum]
\label{def:poincare-covariant-vacuum-hilbert-datum}
Let
\[
  \Poinc=\mathbb R^D\rtimes SO^+(1,D-1)
\]
be the connected Poincare group of \(\mathbb M^D\). A relativistic quantum
theory in the opening framework carries a strongly continuous unitary
representation
\[
  U:\Poinc\to\mathcal U(\Hilb),
  \qquad
  (a,\Lambda)\mapsto U(a,\Lambda),
\]
or the corresponding representation of the double cover when spinorial
representations are present.  The group law is
\[
  (a',\Lambda')(a,\Lambda)
  =
  (a'+\Lambda'a,\Lambda'\Lambda),
\]
so a genuine unitary representation satisfies
\[
  U(a',\Lambda')U(a,\Lambda)
  =
  U(a'+\Lambda'a,\Lambda'\Lambda).
\]
Infinitesimally,
\[
  U(a,\Lambda)
  =
  1+\ii a_\mu P^\mu+\frac{\ii}{2}\omega_{\mu\nu}J^{\mu\nu}
  +O(a^2,\omega^2,a\omega),
\]
where \(P^\mu\) are the self-adjoint energy-momentum generators,
\(P_\mu=\eta_{\mu\nu}P^\nu\), and
\(J^{\mu\nu}=-J^{\nu\mu}\) are the Lorentz generators on a common invariant
domain. The commutation relations are
\[
  [P^\mu,P^\nu]=0,
\]
\[
  [P^\mu,J^{\rho\sigma}]
  =
  -\ii(\eta^{\mu\rho}P^\sigma-\eta^{\mu\sigma}P^\rho),
\]
and
\[
  [J^{\mu\nu},J^{\rho\sigma}]
  =
  -\ii\left(
    \eta^{\nu\rho}J^{\mu\sigma}
    -\eta^{\mu\rho}J^{\nu\sigma}
    -\eta^{\nu\sigma}J^{\mu\rho}
    +\eta^{\mu\sigma}J^{\nu\rho}
  \right).
\]
Equivalently,
\([J^{\rho\sigma},P^\mu]
=\ii(\eta^{\mu\rho}P^\sigma-\eta^{\mu\sigma}P^\rho)\).
The spectrum condition is
\[
  \operatorname{Spec}(P)\subset\overline V_+.
\]
In a vacuum sector there is a unit vector \(\vac\in\Hilb\) satisfying
\[
  U(a,\Lambda)\vac=\vac .
\]
\end{definition}

\section{Translation Spectrum and Mass Sectors}

\begin{proposition}[Joint translation spectrum]
\label{prop:joint-translation-spectrum}
The restriction of the Poincare representation to translations has a unique
joint spectral measure
\[
  E:\{\text{Borel subsets of }\mathbb R^D\}\longrightarrow
  \{\text{orthogonal projections on }\Hilb\}
\]
such that
\[
  P^\mu=\int_{\mathbb R^D}p^\mu\,\dd E(p),
  \qquad
  U(a,\mathbf 1)
  =
  \int_{\mathbb R^D}\ee^{\ii a^\mu p_\mu}\,\dd E(p).
\]
The one-parameter Stone generators \(P^\mu\) are therefore the coordinate
operators of this single joint spectral measure.  In particular the
translation generators are strongly commuting self-adjoint operators: their
spectral projections \(E(\Delta)\) commute.
For a vector \(\psi\in\Hilb\), its energy-momentum spectral measure is
\[
  \mu_\psi(\Delta)=\|E(\Delta)\psi\|^2 ,
  \qquad
  \Delta\subset\mathbb R^D\text{ Borel}.
\]
The spectral support of \(\psi\) is the closed support of this measure.  The
spectrum condition in the preceding section is the statement that
\[
  E(\mathbb R^D\setminus\overline V_+)=0 .
\]
\end{proposition}

\begin{proof}
Restrict \(U\) in
Definition~\ref{def:poincare-covariant-vacuum-hilbert-datum} to the
translation subgroup.  The subgroup is abelian and locally compact, and the
restriction is strongly continuous and unitary.  The Stone--Naimark spectral
theorem therefore gives the single projection-valued measure \(E\) displayed
above.  Explicitly, the equality
\[
  U(a,\mathbf 1)
  =
  \int_{\mathbb R^D}\ee^{\ii a^\mu p_\mu}\,\dd E(p)
\]
is an equality of bounded operators for every \(a\in\mathbb R^D\), and
uniqueness follows because Fourier transforms of finite complex Borel
measures on \(\mathbb R^D\) determine those measures.  Differentiating the
one-parameter subgroup \(s\mapsto U(se_\mu,\mathbf1)\) in Stone's sense gives
the coordinate multiplication operator \(p^\mu\) in this joint spectral
representation, namely
\[
  P^\mu\psi=\int_{\mathbb R^D}p^\mu\,\dd E(p)\,\psi
\]
on its natural domain
\[
  \{\psi\in\Hilb:\int_{\mathbb R^D}|p^\mu|^2\,\dd\|E(p)\psi\|^2<\infty\}.
\]
All spectral projections obtained from \(E\) commute because they are
projections in the same countably additive projection-valued measure.  This
strong commutativity is stronger than the formal infinitesimal relation
\([P^\mu,P^\nu]=0\) on a common invariant domain: for unbounded operators,
commutation on a core does not by itself imply that the spectral projections
commute.  The displayed Lie-algebra commutator is therefore the differential
shadow of the joint functional calculus.  Finally, the operator-valued
support statement
\(E(\mathbb R^D\setminus\overline V_+)=0\) is equivalent to saying that every
scalar spectral measure \(\mu_\psi(\Delta)=\|E(\Delta)\psi\|^2\) is supported
in \(\overline V_+\); this is the precise spectral-measure form of the
spectrum condition.
\end{proof}

\begin{definition}[Invariant mass spectral measure]
\label{def:invariant-mass-spectral-measure}
The invariant mass operator is defined by functional calculus:
\[
  M^2=-P_\mu P^\mu
  =
  \int_{\mathbb R^D}(-p_\mu p^\mu)\,\dd E(p).
\]
For a Borel set \(I\subset[0,\infty)\), the projection onto the corresponding
mass range is
\[
  E_M(I)
  =
  E\bigl(\{p\in\overline V_+:-p_\mu p^\mu\in I\}\bigr).
\]
The vacuum subspace is the zero-momentum subspace
\[
  E(\{0\})\Hilb .
\]
In a pure vacuum sector this subspace is one-dimensional and is spanned by
\(\vac\).
\end{definition}

\begin{definition}[Mass sector and particle shell]
\label{def:mass-sector-and-particle-shell}
  A mass sector of mass \(m\ge0\) is the spectral subspace of \(M^2\) at the
  value \(m^2\), namely \(E_M(\{m^2\})\Hilb\).  For \(m>0\), the corresponding
  positive-energy mass shell is the smooth hyperboloid
  \[
    \Sigma_m^+
    =
    \{p\in\mathbb R^D:p^0=\sqrt{\vec p^{\,2}+m^2},\ p_\mu p^\mu=-m^2\}.
  \]
  For \(m=0\), the nonzero massless particle shell is the punctured future
  light cone
  \[
    \Sigma_0^+
    =
    \{p\in\mathbb R^D:p^0=|\vec p|,\ p_\mu p^\mu=0,\ p\ne0\}.
  \]
  It is isolated when there is an open neighborhood \(N\subset\mathbb R^D\) of
  the shell under discussion such that
  \[
    \operatorname{Spec}(P)\cap N=\Sigma_m^+
  \]
  in the massive case, and analogously with \(\Sigma_0^+\) in a massless
  sector when such a global isolation statement is true.  In a vacuum
  representation the massless cone usually accumulates at the vacuum point
  \(p=0\), so the global isolation condition is a genuine additional spectral
  hypothesis beyond the invariant-mass equation \(M^2=0\).
\end{definition}

\paragraph{Isolated shell as a one-particle subrepresentation.}
An isolated massive particle in a local quantum theory is the closed spectral
subspace
\[
  \Hilb_1=E(\Sigma_m^+)\Hilb
\]
inside the physical Hilbert space.  Its Wigner representation is the
restriction of the Poincare representation to this subspace.  Later, the
K\"all\'en--Lehmann representation detects the same object as an atom in a
two-point spectral measure, and Haag--Ruelle theory uses the isolation
hypothesis to construct asymptotic multi-particle states.
The set \(\Sigma_m^+\) is Borel, so \(E(\Sigma_m^+)\) is an orthogonal
projection.  The spectral covariance relation
\(U(a,\Lambda)E(\Delta)U(a,\Lambda)^{-1}=E(\Lambda\Delta)\) follows from
conjugating the translation representation by \(U(a,\Lambda)\).  Since
\(\Sigma_m^+\) is a positive-energy Lorentz orbit, the range of
\(E(\Sigma_m^+)\) is invariant under the Poincare representation.  Restricting
the representation to this closed range gives the one-particle Wigner
subrepresentation.

\begin{remark}[Infraparticle sectors]
\label{rem:infraparticle-sectors-chapter-two}
The word ``isolated'' is a spectral hypothesis in addition to the mass equation
\(M^2=m^2\).  A theory with a massless long-range gauge field can have charged
sectors whose charged threshold is the lower edge of continuous
energy-momentum support.  In four-dimensional QED, Gauss's law ties nonzero
electric charge to an electromagnetic field at spatial infinity; charged
sector states are therefore represented by operators carrying noncompact
dressing data.  A charged dressing carries a Coulomb tail and admits
arbitrarily soft photon excitations.  The resulting charged particle is an
infraparticle: its spectral support accumulates at the charged threshold in
place of an isolated shell with a nonzero projection \(E(\Sigma_m^+)\).  The
Haag--Ruelle and LSZ constructions below apply to sectors satisfying the
isolated-shell hypothesis; the infrared replacement in QED is treated in
Chapters~\ref{ch:qed-external-states} and
\ref{ch:volume-ii-infrared-inclusive-qed}.
\end{remark}

Figure~\ref{fig:translation-spectrum-mass-shell} displays the spectral
separation assumed in the one-particle subrepresentation construction: an
isolated positive mass shell inside the future cone, separated from the
multiparticle continuum.

\begin{figure}[H]
  \centering
  \resizebox{0.76\textwidth}{!}{%
  \begin{tikzpicture}[x=1cm,y=1cm,every node/.style={font=\scriptsize}]
    \draw[-{Latex[length=2mm]},line width=0.75pt] (-3.4,0)--(3.65,0)
      node[right] {\(\vec p\)};
    \draw[-{Latex[length=2mm]},line width=0.75pt] (0,-0.25)--(0,3.65)
      node[above] {\(p^0\)};
    \draw[gray!55,dashed,line width=0.8pt] (-3.0,3.0)--(0,0)--(3.0,3.0);
    \node[gray!65!black] at (2.72,2.45) {\(\overline V_+\)};
    \draw[blue!70!black,line width=1.2pt]
      plot[smooth,domain=-2.45:2.45,samples=100]
      (\x,{sqrt(\x*\x+0.8)});
    \node[blue!70!black] at (-2.55,2.98) {\(\Sigma_m^+\)};
    \draw[red!75!black,line width=2.0pt]
      plot[smooth,domain=-2.95:2.95,samples=100]
      (\x,{sqrt(\x*\x+2.45)});
    \node[red!75!black,fill=white,inner sep=1pt] at (2.55,3.48)
      {continuum threshold};
    \fill[black] (0,0) circle (1.7pt);
    \node[below right] at (0,0) {\(\vac\)};
    \draw[decorate,decoration={brace,amplitude=5pt},green!45!black]
      (2.03,2.14)--(2.03,2.66);
    \node[green!45!black,align=center,fill=white,inner sep=1pt] at (2.94,2.04)
      {isolating\\neighborhood};
  \end{tikzpicture}
  }
  \caption{The translation spectrum lies in the closed future cone.  A stable
  massive particle corresponds to an isolated positive-energy mass shell,
  separated from the continuum in a neighborhood of the shell.}
  \label{fig:translation-spectrum-mass-shell}
\end{figure}

\section{One-Particle Representations}

\begin{definition}[Massive scalar Wigner one-particle space]
\label{def:massive-scalar-wigner-one-particle-space}
The spectral subspace \(E(\Sigma_m^+)\Hilb\) carries a unitary
representation of the Poincare group.  The corresponding abstract
one-particle representation is described by Wigner's little-group
construction.  For a massive spinless particle, the one-particle Hilbert
space is realized as
\[
  \Hilb_1
  =
  L^2\left(
    \Sigma_m^+,
    \frac{\dd^{D-1}\vec p}{(2\pi)^{D-1}2E_{\vec p}}
  \right),
  \qquad
  E_{\vec p}=\sqrt{\vec p^{\,2}+m^2}.
\]
The measure displayed here is the Lorentz-invariant measure on \(\Sigma_m^+\)
in the chosen normalization. General massive particles carry a
finite-dimensional unitary representation of the double cover of the massive
little group \(SO(D-1)\). Massless finite-helicity particles in four spacetime
dimensions are classified by helicity representations of the massless little
group.
\end{definition}

\begin{definition}[Four-dimensional Wigner orbit datum]
\label{def:four-dimensional-wigner-orbit-datum}
For later reference we record the corresponding four-dimensional Wigner
classification in one place.  Let
\[
  \widetilde{\Poinc}=\R^{1,3}\rtimes SL(2,\C)
\]
be the universal-cover version of the connected Poincare group.  An
irreducible unitary representation whose translation spectrum is supported on
a single Lorentz orbit is obtained by choosing a representative momentum
\(\bar p\), an irreducible unitary representation of its stabilizer
\[
  G_{\bar p}=\{\Lambda\in SL(2,\C):\Lambda \bar p=\bar p\},
\]
and inducing to \(\widetilde{\Poinc}\).  The spectrum condition keeps only
orbits contained in the closed future cone, but the representation-theoretic
list is useful because it makes explicit which cases are being assumed away.
\end{definition}

\begin{table}[H]
  \centering
  \scriptsize
  \renewcommand{\arraystretch}{1.25}
  \setlength{\tabcolsep}{2pt}
  \begin{tabular}{|
    >{\raggedright\arraybackslash}p{0.14\textwidth}|
    >{\raggedright\arraybackslash}p{0.17\textwidth}|
    >{\raggedright\arraybackslash}p{0.14\textwidth}|
    >{\raggedright\arraybackslash}p{0.19\textwidth}|
    >{\raggedright\arraybackslash}p{0.18\textwidth}|}
    \hline
    Orbit type & Representative and invariant & Little group & Irreducible
    little-group data & Role in this monograph \\
    \hline
    Vacuum orbit &
    \(\bar p=0\) &
    \(SL(2,\C)\) &
    The pure vacuum sector uses the trivial representation. &
    This is the zero-momentum state, not a particle shell. \\
    \hline
    Massive positive-energy orbit &
    \(\bar p=(m,0,0,0)\), \(m>0\), \(\bar p^2=-m^2\) &
    \(SU(2)\) &
    Spin \(j=0,\frac12,1,\frac32,\ldots\), with
    \(\dim V_j=2j+1\). &
    Isolated massive particles and their higher-spin tower; developed in
    Chapters~\ref{ch:massive-particles-spin} and
    \ref{ch:spinor-fields-fermionic-path-integrals}. \\
    \hline
    Massless finite-helicity orbit &
    \(\bar p=(\kappa,0,0,\kappa)\), \(\kappa>0\), \(\bar p^2=0\) &
    double cover of \(ISO(2)\) &
    The two null-translation generators act trivially; the rotation generator
    has helicity \(h\in\frac12\mathbb Z\). &
    Photons, gravitons, and finite-helicity external states; developed in
    Chapters~\ref{ch:massless-particles-helicity-gauge-redundancy} and
    \ref{ch:maxwell-theory-constraints-gauge-fixing}. \\
    \hline
    Massless continuous-spin orbit &
    Same lightlike orbit, \(\bar p^2=0\), with nontrivial little-group
    translations &
    double cover of \(ISO(2)\) &
    Nonzero transverse-translation orbit
    \(N_1^2+N_2^2=\rho^2>0\), giving an infinite helicity tower mixed by
    boosts. &
    A legitimate positive-energy Wigner class, but not generated by
    point-local covariant fields under the usual locality assumptions; see
    Remark~\ref{rem:continuous-spin-scope}. \\
    \hline
    Spacelike, or tachyonic, orbit &
    \(\bar p=(0,\mu,0,0)\), \(\mu>0\), \(\bar p^2=\mu^2\) &
    double cover of \(SO^+(1,2)\) &
    Unitary little-group representations are representations of a noncompact
    group, typically infinite-dimensional. &
    Excluded here by the spectrum condition: the full Lorentz orbit is not
    contained in \(\overline V_+\) and has no invariant lower energy bound. \\
    \hline
  \end{tabular}
  \caption{Four-dimensional Wigner orbit classes for irreducible
  one-particle representations of the universal cover of the connected
  Poincare group.  A general Hilbert space representation may contain direct
  sums or direct integrals of such irreducibles; a stable particle species is
  an isolated copy of one of the particle orbits satisfying the spectral and
  localization hypotheses used later.}
  \label{tab:wigner-classification-four-dimensional}
\end{table}

This construction supplies one-particle state spaces and their Poincare
transformation law. A local quantum theory also supplies a localization
assignment for observables.

\section{Free Fock Space From One-Particle Data}

\begin{definition}[Free bosonic and fermionic Fock spaces]
\label{def:free-bosonic-fermionic-fock-spaces}
From a one-particle Hilbert space \(\Hilb_1\), define the bosonic Fock space
\[
  \mathcal F_s(\Hilb_1)
  =
  \bigoplus_{n=0}^{\infty}\operatorname{Sym}^n\Hilb_1
\]
and the fermionic Fock space
\[
  \mathcal F_a(\Hilb_1)
  =
  \bigoplus_{n=0}^{\infty}\wedge^n\Hilb_1.
\]
\end{definition}

\begin{definition}[Bosonic creation-operator tensor normalization]
\label{def:bosonic-creation-operator-tensor-normalization}
The symmetric tensor powers use the creation-operator normalization.  Let
\[
  \Pi_{s,n}
  =
  \frac{1}{n!}
  \sum_{\sigma\in S_n}U_\sigma
\]
be the orthogonal projection from \(\Hilb_1^{\otimes n}\) to its symmetric
subspace, where
\[
  U_\sigma(\psi_1\otimes\cdots\otimes\psi_n)
  =
  \psi_{\sigma^{-1}(1)}\otimes\cdots\otimes\psi_{\sigma^{-1}(n)} .
\]
For \(\psi_1,\ldots,\psi_n\in\Hilb_1\) define
\begin{equation}
  \psi_1\odot\cdots\odot\psi_n
  :=
  \sqrt{n!}\,
  \Pi_{s,n}(\psi_1\otimes\cdots\otimes\psi_n)
  \in \operatorname{Sym}^n\Hilb_1 .
  \label{eq:bosonic-fock-odot-normalization}
\end{equation}
\end{definition}

\paragraph{Bosonic Fock inner product in creation-operator normalization.}
With this convention
\begin{equation}
  \bigl\langle
  \psi_1\odot\cdots\odot\psi_n,\,
  \chi_1\odot\cdots\odot\chi_m
  \bigr\rangle
  =
  \delta_{mn}
  \sum_{\sigma\in S_n}
  \prod_{j=1}^n
  \langle\psi_j,\chi_{\sigma(j)}\rangle_{\Hilb_1}.
  \label{eq:bosonic-fock-odot-inner-product}
\end{equation}
Different tensor powers in the direct sum defining
\(\mathcal F_s(\Hilb_1)\) are orthogonal, giving the factor \(\delta_{mn}\).
For \(m=n\), Definition~\ref{def:bosonic-creation-operator-tensor-normalization}
gives
\[
  \psi_1\odot\cdots\odot\psi_n
  =
  \frac1{\sqrt{n!}}\sum_{\sigma\in S_n}
  \psi_{\sigma^{-1}(1)}\otimes\cdots\otimes\psi_{\sigma^{-1}(n)}
\]
and the analogous formula for the \(\chi\)'s.  Hence
\[
  \bigl\langle
  \psi_1\odot\cdots\odot\psi_n,\,
  \chi_1\odot\cdots\odot\chi_n
  \bigr\rangle
  =
  \frac1{n!}\sum_{\sigma,\tau\in S_n}
  \prod_{i=1}^n
  \langle\psi_{\sigma^{-1}(i)},\chi_{\tau^{-1}(i)}\rangle_{\Hilb_1}.
\]
For each fixed \(\sigma\), set \(j=\sigma^{-1}(i)\) and
\(\rho=\tau^{-1}\sigma\).  The product becomes
\(\prod_{j=1}^n\langle\psi_j,\chi_{\rho(j)}\rangle_{\Hilb_1}\), and as
\(\tau\) runs over \(S_n\), so does \(\rho\).  The sum over \(\sigma\) cancels
the prefactor \(1/n!\), leaving the displayed permutation sum.

\begin{definition}[Mass-shell Fock test domain and kernels]
\label{def:mass-shell-fock-test-domain-kernels}
For the massive scalar one-particle space
\(\Hilb_1=L^2(\Sigma_m^+,\dd\mu_m)\), identify the mass shell with
\(\mathbb R^{D-1}\) by
\[
  \vec p\longmapsto (E_{\vec p},\vec p),
  \qquad E_{\vec p}=\sqrt{\vec p^{\,2}+m^2},
\]
and define
\[
  \mathcal S_m=\mathcal S(\Sigma_m^+)
\]
to be the wave packets whose pullback to the \(\vec p\)-chart is a Schwartz
function.  The space \(\mathcal S_m\) is dense in \(\Hilb_1\), is contained in
the domain of every polynomial in the one-particle energy-momentum operators,
and is invariant under the one-particle Poincare representation.  The common
algebraic Fock domain used below is
\[
  \mathcal D_0
  =
  \operatorname{span}\{
    a^\dagger(\psi_1)\cdots a^\dagger(\psi_n)\vac
    \mid
    n\ge0,\ \psi_1,\ldots,\psi_n\in\mathcal S_m
  \}.
\]
It is dense in \(\mathcal F_s(\Hilb_1)\) and invariant under creation and
annihilation operators smeared by elements of \(\mathcal S_m\).
For a free scalar species, the momentum-space annihilation and creation
operators \(a(\vec p)\) and \(a^\dagger(\vec p)\) are operator-valued
distributions characterized by
\[
  [a(\vec p),a^\dagger(\vec q)]
  =
  (2\pi)^{D-1}2E_{\vec p}\,\delta^{(D-1)}(\vec p-\vec q),
\]
with all other commutators vanishing. The free Hamiltonian is
\[
  H_0
  =
  \int
  \frac{\dd^{D-1}\vec p}{(2\pi)^{D-1}2E_{\vec p}}\,
  E_{\vec p}\,a^\dagger(\vec p)a(\vec p),
\]
after the vacuum energy has been normalized.
The spatial momentum is
\[
  \vec P_0
  =
  \int
  \frac{\dd^{D-1}\vec p}{(2\pi)^{D-1}2E_{\vec p}}\,
  \vec p\,a^\dagger(\vec p)a(\vec p).
\]
Write
\[
  \dd\mu_m(\vec p)
  =
  \frac{\dd^{D-1}\vec p}{(2\pi)^{D-1}2E_{\vec p}} .
\]
For a one-particle test wave packet \(\psi\in\mathcal S_m\), the smeared
creation and annihilation operators on \(\mathcal D_0\) are
\[
  a^\dagger(\psi)
  =
  \int \dd\mu_m(\vec p)\,\psi(\vec p)a^\dagger(\vec p),
  \qquad
  a(\psi)
  =
  \int \dd\mu_m(\vec p)\,\overline{\psi(\vec p)}a(\vec p).
\]
The closed Fock operators can be defined for all \(\psi\in\Hilb_1\).  The
operator-valued distribution notation and the common-domain commutator
calculation use the dense test class \(\mathcal S_m\).  For
\(\psi,\chi\in\mathcal S_m\), the operators on \(\mathcal D_0\) satisfy
\[
  [a(\psi),a^\dagger(\chi)]
  =
  \int \dd\mu_m(\vec p)\,\overline{\psi(\vec p)}\chi(\vec p)
  =
  \langle\psi,\chi\rangle_{\Hilb_1}.
\]
\end{definition}

\begin{definition}[Noncovariant creation and momentum-ket normalization]
\label{def:noncovariant-creation-momentum-ket-normalization}
There is a second common normalization, useful for displaying the elementary
Fock construction.  Define operator-valued distributions
\(\mathfrak a_{\vec p}\) by
\[
  a(\vec p)
  =
  \sqrt{(2\pi)^{D-1}2E_{\vec p}}\;\mathfrak a_{\vec p}.
\]
Then
\[
  [\mathfrak a_{\vec p},\mathfrak a_{\vec q}^{\dagger}]
  =
  \delta^{(D-1)}(\vec p-\vec q),
  \qquad
  [\mathfrak a_{\vec p},\mathfrak a_{\vec q}]=
  [\mathfrak a_{\vec p}^{\dagger},\mathfrak a_{\vec q}^{\dagger}]=0.
\]
The corresponding sharp-momentum kets are distributional vectors,
\[
  |\vec p\rangle=\mathfrak a_{\vec p}^{\dagger}\vac,
  \qquad
  |\vec p_1,\ldots,\vec p_n\rangle
  =
  \mathfrak a_{\vec p_1}^{\dagger}\cdots
  \mathfrak a_{\vec p_n}^{\dagger}\vac,
\]
and
\[
  \langle\vec p|\vec q\rangle=\delta^{(D-1)}(\vec p-\vec q).
\]
In this normalization
\[
  H_0=\int\dd^{D-1}\vec p\,
  E_{\vec p}\,\mathfrak a_{\vec p}^{\dagger}\mathfrak a_{\vec p},
  \qquad
  \vec P_0=\int\dd^{D-1}\vec p\,
  \vec p\,\mathfrak a_{\vec p}^{\dagger}\mathfrak a_{\vec p}.
\]
\end{definition}

\begin{definition}[Free multi-species Fock representation]
\label{def:free-multi-species-fock-representation}
For a finite collection of free stable species, let \(B\) be the set of
bosonic species and \(F\) the set of fermionic species, with one-particle
spaces \(\Hilb_{1,b}\) and \(\Hilb_{1,f}\).  The free multi-species Hilbert
space is the completed tensor product
\[
  \Hilb_{\rm Fock}
  =
  \bigotimes_{b\in B}\mathcal F_s(\Hilb_{1,b})
  \otimes
  \bigotimes_{f\in F}\mathcal F_a(\Hilb_{1,f}) .
\]
The vacuum is the tensor product of the factor vacua, and the free Poincare
representation is the tensor product of the corresponding second-quantized
one-particle representations.  Equivalently,
\(\bigotimes_{b\in B}\mathcal F_s(\Hilb_{1,b})\simeq
\mathcal F_s(\bigoplus_{b\in B}\Hilb_{1,b})\), and similarly for fermions with
\(\mathcal F_a\).  The operator algebra on the fermionic factors uses the
graded tensor-product convention, so homogeneous odd creation or annihilation
operators associated with distinct fermion species anticommute.
\end{definition}

The Fock construction gives a free multi-particle representation. A local QFT
also specifies which operators are localized in each bounded spacetime region.

\section{Locality As A Constraint On Interacting Particle Dynamics}

At a cutoff or algebraic bookkeeping level, an interaction written in free
Fock variables has the form
\[
  H=H_0+H_{\mathrm{int}},
\]
where \(H_{\mathrm{int}}\) is a sum of normally ordered monomials such as
\[
  \int K_{2,1}\,
  a^\dagger a^\dagger a
  +
  \int K_{1,2}\,
  a^\dagger a a
  +
  \cdots ,
\]
with kernels \(K_{r,s}\) and the mass-shell measures suppressed in the
displayed shorthand.  This expression is not yet a definition of a
relativistic local theory.  A local-theory datum also includes translation
generators \(P_\mu\), Lorentz generators \(J^{\mu\nu}\), a Poincare-invariant
vacuum, and local observables or fields satisfying covariance and
microcausality.  In an interacting theory the boost generators, the
Hamiltonian, and the local field operators are tied together by the Poincare
algebra and locality.

The causal input is spacetime-local operation.  A local operation near a point
or region can influence only its causal future.  The algebraic expression of
this statement is locality for spacelike separated regions.

Figure~\ref{fig:causal-lightcone-local-operation} fixes the light-cone convention
used here: the local operation sits at the vertex, the allowed influence lies in
the causal future, and spacelike separated regions remain outside this causal
domain.

\begin{figure}[H]
  \centering
  \resizebox{0.72\textwidth}{!}{%
  \begin{tikzpicture}[scale=0.95,>=Stealth]
    \draw[->,line width=0.75pt] (-0.1,-2.45) -- (-0.1,2.65)
      node[above] {\(x^0\)};
    \draw[->,line width=0.75pt] (-2.65,0) -- (3.45,0)
      node[right] {\(\vec x\)};
    \draw[green!45!black,line width=1.15pt] (0,0)--(-1.75,2.25);
    \draw[green!45!black,line width=1.15pt] (0,0)--(1.75,2.25);
    \draw[green!45!black,line width=1.15pt] (0,0)--(-1.75,-2.25);
    \draw[green!45!black,line width=1.15pt] (0,0)--(1.75,-2.25);
    \node[green!45!black] at (2.25,2.38) {light cone};
    \fill[blue!75!black] (0,0) circle (2pt);
    \node[blue!75!black,below left] at (0,0) {local operation};
    \draw[purple,dashed,line width=0.95pt,->]
      (0,0).. controls (0.25,0.55) and (0.35,1.2)..(0.62,1.92);
    \fill[purple] (0.62,1.92) circle (2.8pt);
    \draw[red!78!black,dashed,line width=0.95pt,->] (0,0)--(2.15,0.92);
    \fill[red!78!black] (2.15,0.92) circle (2.8pt);
    \node[red!78!black,align=left,anchor=west] at (2.35,0.92)
      {spacelike signal\\excluded by locality};
  \end{tikzpicture}
  }
  \caption{Causal localization in Minkowski space.  The field-theoretic
  condition corresponding to the exclusion of spacelike signaling is
  commutativity, or graded commutativity, of observables localized in
  spacelike separated regions.}
  \label{fig:causal-lightcone-local-operation}
\end{figure}

\section{Locality Data}

\begin{definition}[Local observable assignment]
\label{def:local-observable-assignment-chapter-two}
For bounded open regions \(\mathcal O\subset\mathbb M^D\), the opening
framework assigns algebras \(\Obs(\mathcal O)\subset\mathcal B(\Hilb)\). If
\(\mathcal O_1\) and \(\mathcal O_2\) are spacelike separated, locality is
\[
  [A,B]=0,
  \qquad
  A\in\Obs(\mathcal O_1),\quad B\in\Obs(\mathcal O_2),
\]
with the graded commutator after a fermion parity grading has been specified.

This condition states compatibility of spacelike separated local observables:
their order gives the same matrix elements in every vector for which the
matrix elements are defined.
\end{definition}

\section{Free Scalar Field As A Derived Example}

\begin{definition}[Free scalar field from mass-shell creation kernels]
\label{def:free-scalar-field-mass-shell-kernels}
Use the dense domain \(\mathcal D_0\subset\mathcal F_s(\Hilb_1)\) generated
from the mass-shell Schwartz test space \(\mathcal S_m\).  On \(\mathcal D_0\),
define the free scalar field as the
operator-valued distribution
\[
  \widehat\phi(x)
  =
  \int
  \frac{\dd^{D-1}\vec p}{(2\pi)^{D-1}2E_{\vec p}}
  \left(
    a(\vec p)\ee^{\ii p\cdot x}
    +
    a^\dagger(\vec p)\ee^{-\ii p\cdot x}
  \right),
\]
where \(p^0=E_{\vec p}\) and
\[
  p\cdot x=\eta_{\mu\nu}p^\mu x^\nu.
\]
For \(f\in C_c^\infty(\mathbb M^D)\), the smeared field is
\[
  \widehat\phi(f)=\int\dd^Dx\,f(x)\widehat\phi(x).
\]
In the noncovariant normalization this is equivalently
\[
  \widehat\phi(x)
  =
  \int
  \frac{\dd^{D-1}\vec p}
       {\sqrt{(2\pi)^{D-1}2E_{\vec p}}}
  \left(
    \mathfrak a_{\vec p}\ee^{\ii p\cdot x}
    +
    \mathfrak a_{\vec p}^{\dagger}\ee^{-\ii p\cdot x}
  \right).
\]
\end{definition}

\paragraph{Covariance of the free scalar field.}
The free scalar field is covariant.  Let \(g=(b,\Lambda)\in \Poinc\), and let
\(U_1(g)\) be the one-particle scalar representation on \(\Hilb_1\).  The
free Fock representation is its second quantization \(U_0(g)=\Gamma(U_1(g))\),
so
\[
  U_0(g)a^\dagger(\psi)U_0(g)^{-1}
  =
  a^\dagger(U_1(g)\psi)
\]
for \(\psi\in\mathcal S_m\) on \(\mathcal D_0\).  Since \(U_1(g)\mathcal S_m
\subset\mathcal S_m\), this identity preserves the common test domain; the
closed operators give the corresponding Hilbert-space identity for arbitrary
\(\psi\in\Hilb_1\).  The adjoint identity gives the annihilation part.
Applying these identities to the distributional creation and annihilation
kernels in Definition~\ref{def:free-scalar-field-mass-shell-kernels}, followed
by the invariant change of variables on \(\Sigma_m^+\), gives
\[
  U_0(b,\Lambda)\widehat\phi(x)U_0(b,\Lambda)^{-1}
  =
  \widehat\phi(\Lambda x+b).
\]
Thus the calculation checks compatibility of the free kernel definition with
the chosen one-particle representation.

\begin{proposition}[Free scalar microcausality]
\label{prop:free-scalar-microcausality-chapter-two}
The commutator is
\[
  [\widehat\phi(x),\widehat\phi(y)]
  =
  \ii\Delta(x-y),
\]
where the Pauli--Jordan distribution is
\[
  \Delta(x)
  =
  \frac{1}{\ii}
  \int
  \frac{\dd^{D-1}\vec p}{(2\pi)^{D-1}2E_{\vec p}}
  \left(
    \ee^{\ii p\cdot x}-\ee^{-\ii p\cdot x}
  \right).
\]
The distribution \(\Delta\) is Lorentz invariant and its restriction to the
spacelike region vanishes. Consequently
\[
  [\widehat\phi(f),\widehat\phi(g)]=0
\]
when the supports of \(f\) and \(g\) are spacelike separated.
\end{proposition}

\begin{proof}
The canonical commutation relation in
Definition~\ref{def:mass-shell-fock-test-domain-kernels} gives the displayed
Pauli--Jordan commutator on the common domain \(\mathcal D_0\).  In smeared
form, for \(f,g\in C_c^\infty(\mathbb M^D)\),
\[
  [\widehat\phi(f),\widehat\phi(g)]
  =
  \ii\int \dd^Dx\,\dd^Dy\, f(x)g(y)\Delta(x-y)
\]
as a scalar multiple of the identity on \(\mathcal D_0\).  The mass-shell
measure and the two exponentials are invariant under the connected Lorentz
group, so \(\Delta(\Lambda x)=\Delta(x)\) as a distribution.  On the open
spacelike region \(x^2>0\), the mass-shell Fourier integral defines a smooth
solution of the Klein--Gordon equation.  For each spacelike \(x\) choose a
proper orthochronous Lorentz transformation sending it to \((0,\vec r)\).
At equal time the two exponentials in the defining integral are exchanged by
the measure-preserving map \(\vec p\mapsto-\vec p\), hence the smooth
restriction of \(\Delta\) to that point is zero.  Since the spacelike region
is open and every point in it is covered by this argument, \(\Delta\) vanishes
there as a distribution.  If the supports of \(f\) and \(g\) are spacelike
separated, then \(x-y\) remains in this open spacelike region on
\(\operatorname{supp}f\times\operatorname{supp}g\); the displayed smeared
commutator therefore vanishes.
\end{proof}

\begin{remark}
\label{rem:scalar-odd-locality-fails}
The cancellation above uses the antisymmetric combination
\(\Delta_+(x-y)-\Delta_+(y-x)\), equivalently the ordinary commutator for an
even scalar field.  If the same scalar two-point function were assigned odd
locality, the vacuum matrix element of the required anticommutator would contain
\[
  \Delta_+(x-y)+\Delta_+(y-x).
\]
At equal time and spacelike separation this becomes
\[
  2\int
  \frac{\dd^{D-1}\vec p}{(2\pi)^{D-1}2E_{\vec p}}\,
  \cos(\vec p\cdot\vec r),
\]
which has spacelike support.  Thus the sign in the Pauli--Jordan distribution
is already the local free-field manifestation of the spin-statistics
connection.  The theorem-level statement, with its positivity, covariance, and
domain hypotheses, is discussed in
Chapter~\ref{ch:spinor-fields-fermionic-path-integrals}.
\end{remark}

This calculation establishes locality for the free scalar field from the
stated data: the free scalar Fock representation, the Lorentz-invariant
mass-shell measure, the distributional field definition, and smearing by test
functions.

\section{Constructed Structures}

The chapter has constructed five distinct structures. Hilbert-space quantum
mechanics supplies states and operators. Translation spectral theory supplies
energy-momentum support, invariant mass, and particle subspaces inside the
physical Hilbert space. Poincare representation theory supplies the
transformation law of one-particle state spaces. Fock space supplies free
multi-particle state spaces. A local quantum theory supplies local observable
assignments satisfying covariance and microcausality. The free scalar field
realizes these structures explicitly and will serve as the first model for
the general field framework.
